
\paragraph{Soft margin classification}

\begin{equation}\label{eq:soft-margin-classification}
  \min_{w \in \R^d, \xi \in \R^M, b \in \R} \frac{1}{2} \norm*{w}_2^2 + C \sum_{i=1}^M \xi_i
  \quad \text{s.t.} \quad
  \begin{cases}
    y_i\left(\inner{w, x_i} + b\right) \geq 1 - \xi_i, & i = 1, \ldots, M \\
    \xi_i \geq 0,                                      & i = 1, \ldots, M
  \end{cases} \tag{3}
\end{equation}
\paragraph{The dual problem}
\begin{equation}
  Y = \operatorname{diag}(y_1, \ldots, y_M) \in \R^{M \times M}, \quad G = \left(\inner{x_i, x_j}\right)_{i,j=1}^M \in \R^{M \times M}, \quad \alpha = (\alpha_1, \ldots, \alpha_M)^\top \in \R^M
\end{equation}

where $Y$ is a diagonal matrix with the labels $y_i$ on the diagonal, and $G$ is the Gram matrix of the data points $x_i$. The dual problem can be written as

\begin{equation}
  \min_{\alpha \in \R^M} \frac{1}{2} \inner{\alpha, YGY\alpha} - \inner{\mathbf{1}_M, \alpha}
  \quad \text{s.t.} \quad
  \begin{cases}
    \inner{y, \alpha} = 0                      \\
    0 \leq \alpha_i \leq C, & i = 1, \ldots, M
  \end{cases} \tag{7}
\end{equation}

Discuss existence and uniqueness of equation (3) and (7).

\paragraph{Existence of the primal problem}

\begin{lemma}{Feasibility of the Primal Problem}{}
    The primal problem is feasible, i.e. the feasible set is non-empty, if there exists at least one point $(\symbf{w}, b, \symbf{\xi})$ such that the constraints are satisfied.
\end{lemma}

\begin{proof}{}{}
  We show that the primal problem is feasible by constructing a solution $(\symbf{w}, b, \symbf{\xi})$ that satisfies all constraints.

  Let $\symbf{w} = \symbf{0}$, $b = 0$, and $\symbf{\xi}_i = 1$ for all $i = 1, \ldots, M$. Then:

  \[
    y_i\left(\inner*{\symbf{w}, \symbf{x}_i} + b\right) = y_i\left(\inner*{\symbf{0}, \symbf{x}_i} + 0\right) = 0
  \]

  Since $0 \geq 1 - \symbf{\xi}_i$ when $\symbf{\xi}_i \geq 1$, and we've set $\symbf{\xi}_i = 1$, the first constraint is satisfied:

  \[
    y_i\left(\inner*{w, x_i} + b\right) \geq 1 - \xi_i, \quad i = 1, \ldots, M
  \]

  The second constraint $\xi_i \geq 0$ for all $i = 1, \ldots, M$ is clearly satisfied since $\xi_i = 1 > 0$.

  Therefore, $(\symbf{0}, 0, \symbf{1})$ is a feasible solution, which proves that the primal problem is feasible.

\end{proof}

\begin{lemma}{Boundedness of the Primal Problem}{}
  Since $C > 0$, $\symbf{\xi}_i \geq 0$, and $\norm*{\symbf{w}}_2^2 \geq 0$, the objective function is trivially bounded from below by $0$.
    \[
        \norm*{\symbf{w}}_2^2 + C \sum_{i=1}^M \xi_i \geq 0
    \]
\end{lemma}
\begin{lemma}{Coercitivity of the Primal Problem}{}
  The primal problem is coercive, i.e. the objective function tends to infinity as $\norm*{\symbf{w}}_2^2 + C \sum_{i=1}^M \xi_i \to \infty$.
    \[
        \lim_{\norm*{\symbf{w}}_2^2 + C \sum_{i=1}^M \xi_i \to \infty} \left(\frac{1}{2} \norm*{\symbf{w}}_2^2 + C \sum_{i=1}^M \xi_i\right) = \infty
    \]
\end{lemma}




